\chapter{Lagrange Multipliers}\index{Lagrange multipliers}

Imagine hiking on a mountain, but you must stay on a marked trail. The highest
point you can reach is probably not the summit. It is the highest point
\emph{on the trail}. Most real optimization problems look like this. A company
wants the most profit, but it has a fixed budget. An engineer wants the strongest
beam, but it must fit in a given space. A can maker wants to use the least
aluminum, but the can must hold 355 cm$^3$.

In the Optimization chapter, you solved problems like these by using the
constraint to eliminate a variable. That works when the algebra cooperates.
This chapter gives a method that works even when it does not, and it comes with
a beautiful picture: at the best point, two gradients line up.

\section{Constrained Optimization}\index{constrained optimization}
The general problem is:
$$\text{maximize (or minimize) } f(x, y) \quad \text{subject to} \quad g(x, y) = c$$
The function $f$ is the \newterm{objective}: the thing you care about. The equation
$g(x, y) = c$ is the \newterm{constraint}: a curve in the plane that you must stay on.

Here is a simple example to keep in mind. You have 16 m of fence to make a
rectangle. If the sides are $x$ and $y$, then $2x + 2y = 16$, or $x + y = 8$, and you
want to maximize the area $f(x, y) = xy$. (You can solve this one by substitution,
but it is a good test case for the new method.)

\section{The Picture}
Draw the level curves of $f(x, y) = xy$: the hyperbolas $xy = 4$, $xy = 9$, $xy = 16$,
and $xy = 25$ (the Multivariate Functions chapter introduced level curves). Then
draw the constraint, the line $x + y = 8$ (figure \ref{fig:lagrange-contours}).

\begin{figure}[H]
\centering
\begin{tikzpicture}[scale=0.75]
\draw[-latex] (0,0) -- (9.4,0) node[below right] {$x$};
\draw[-latex] (0,0) -- (0,9.2) node[above] {$y$};
\foreach \x in {2,4,6,8} { \draw (\x,0.1) -- (\x,-0.1) node[below] {\small \x}; \draw (0.1,\x) -- (-0.1,\x) node[left] {\small \x}; }
\foreach \c/\lo in {4/0.45, 9/1.0, 16/1.8, 25/2.8} {
  \draw[gray, thick, domain=\lo:9, samples=80, smooth] plot (\x, {\c/\x});
  \node[gray, right] at (9, {\c/9}) {\small $xy = \c$};
}
\draw[very thick, sdkblue] (0,8) -- (8,0);
\node[sdkblue, right] at (0.3,8.3) {constraint $x + y = 8$};
\fill[orange] (1.354,6.646) circle (0.12);
\fill[orange] (6.646,1.354) circle (0.12);
\fill[red] (4,4) circle (0.14) node[above right] {$(4, 4)$};
\end{tikzpicture}
\caption{Level curves of the area $f = xy$ (gray) and the fence constraint
$x + y = 8$ (blue). The line crosses $xy = 9$ at the orange points but only
touches $xy = 16$, at $(4, 4)$.}
\label{fig:lagrange-contours}
\end{figure}

Walk along the blue line from the top left. At first, you cross level curves with
bigger and bigger values: $xy = 4$, then $xy = 9$. You are going uphill. At $(4, 4)$,
the line just \emph{touches} the level curve $xy = 16$ without crossing it, and after
that you start crossing the same curves again in reverse: downhill. The curve
$xy = 25$ is out of reach entirely. So the best area is 16 m$^2$, at $x = y = 4$.

That is the whole idea: \textbf{at a constrained maximum or minimum, the constraint
curve is tangent to a level curve of $f$.} If the path crossed the level curve
instead, you could keep walking in one direction and climb higher.

\subsection{Tangent Curves Have Parallel Gradients}
From the Gradients chapter, the gradient $\nabla f$ points in the direction of
steepest ascent, and it is perpendicular to the level curve through that point
(moving along a level curve does not change $f$, so the gradient has no component
along it). The constraint curve is a level curve of $g$, so $\nabla g$ is
perpendicular to the constraint.

When the two curves are tangent, they share a tangent line, so their perpendiculars
point along the same line. That means the gradients are parallel: one is a multiple
of the other.

\begin{figure}[H]
\centering
\begin{tikzpicture}[scale=0.95]
\begin{scope}
\clip (0.4,0.4) rectangle (8.6,8.0);
\draw[gray, thick, domain=1.9:8.4, samples=80, smooth] plot (\x, {16/\x});
\draw[gray, thick, domain=1.1:8.4, samples=80, smooth] plot (\x, {9/\x});
\draw[very thick, sdkblue] (0.6,7.4) -- (7.4,0.6);
\end{scope}
\node[gray, right] at (8.4, 16/8.4) {\small $xy = 16$};
\node[gray, right] at (8.4, 9/8.4) {\small $xy = 9$};
\node[sdkblue, below left] at (7.4,0.6) {constraint};
% at the optimum (4,4): grad f = (4,4), grad g = (1,1)
\draw[-latex, very thick, red] (4,4) -- (5.0,5.0) node[above right] {$\nabla f$};
\draw[-latex, thick, sdkblue] (4,4) -- (4.6,4.6);
\node[sdkblue, right] at (4.6,4.45) {$\nabla g$};
\fill[red] (4,4) circle (0.09);
\node[red, below left] at (4,4) {optimum};
% at a crossing point (1.354, 6.646): grad f = (6.646, 1.354), scaled by 1/4
\draw[-latex, very thick, red] (1.354,6.646) -- (3.016,6.985) node[right] {$\nabla f$};
\draw[-latex, thick, sdkblue] (1.354,6.646) -- (1.954,7.246) node[left, xshift=-2pt] {$\nabla g$};
\draw[-latex, very thick, orange] (1.30,6.42) -- (2.10,5.62);
\node[orange, right] (up) at (4.4,6.2) {\small uphill along the path};
\draw[orange, thin] (up.west) -- (1.78,5.98);
\fill[orange] (1.354,6.646) circle (0.09);
\end{tikzpicture}
\caption{At the optimum, $\nabla f$ and $\nabla g$ point the same way. At the orange
point, $\nabla f$ tilts away from $\nabla g$. The part of $\nabla f$ that runs along the
blue path (orange arrow) says you can still increase $f$ by moving along the path.}
\label{fig:lagrange-gradients}
\end{figure}

Figure \ref{fig:lagrange-gradients} shows the difference. At the orange point,
$\nabla f = (6.65, 1.35)$ is not parallel to $\nabla g = (1, 1)$. Part of $\nabla f$
points along the fence line, so sliding along the line raises the area. At $(4, 4)$,
$\nabla f = (4, 4)$ is exactly $4\nabla g$. No part of it points along the line, so
there is nowhere better to slide.

\section{The Method}
Write the parallel-gradient condition with a multiplier $\lambda$ (the Greek letter
lambda), called a \newterm{Lagrange multiplier}, after Joseph-Louis Lagrange, who
invented the method in 1788.

\begin{mdframed}[style = important, frametitle = {Lagrange Multipliers}]
To find the maximum and minimum of $f(x, y)$ subject to $g(x, y) = c$, solve
$$\nabla f = \lambda\nabla g \qquad \text{and} \qquad g(x, y) = c$$
that is, the three equations
$$f_x = \lambda g_x, \qquad f_y = \lambda g_y, \qquad g(x, y) = c$$
for $x$, $y$, and $\lambda$. Then evaluate $f$ at every solution. The largest value is
the maximum, and the smallest is the minimum.
\end{mdframed}

The method finds \emph{candidates}, just as setting a derivative to zero did in one
variable. Always compare the values of $f$ at all of them.

\textbf{Example}: Maximize $f(x, y) = xy$ subject to $x + y = 8$.

\textbf{Solution}: $\nabla f = (y, x)$ and $\nabla g = (1, 1)$. The equations are
$$y = \lambda, \qquad x = \lambda, \qquad x + y = 8$$
So $x = y$, and then $x = y = 4$, with $\lambda = 4$. The maximum area is $16$, just as
the picture showed.

\textbf{Example}: Find the maximum and minimum of $f(x, y) = x^2 + 2y^2$ on the unit
circle $x^2 + y^2 = 1$.

\textbf{Solution}: $\nabla f = (2x, 4y)$ and $\nabla g = (2x, 2y)$. The equations are
$$2x = 2\lambda x, \qquad 4y = 2\lambda y, \qquad x^2 + y^2 = 1$$
The first says $x(1 - \lambda) = 0$, so $x = 0$ or $\lambda = 1$. The second says
$y(2 - \lambda) = 0$, so $y = 0$ or $\lambda = 2$.
\begin{itemize}
\item If $\lambda = 1$, then $y = 0$, and the constraint gives $x = \pm 1$. There, $f = 1$.
\item If $\lambda = 2$, then $x = 0$, and the constraint gives $y = \pm 1$. There, $f = 2$.
\end{itemize}
($x = y = 0$ is not on the circle.) The maximum is 2, at $(0, \pm 1)$, and the minimum
is 1, at $(\pm 1, 0)$. Notice that one set of equations found both.

\section{The Soda Can, Without Substitution}
In the Optimization chapter, you found the can that holds 355 cm$^3$ using the least
aluminum. Here is the same problem with Lagrange multipliers. Minimize the surface
area $A(r, h) = 2\pi r^2 + 2\pi rh$ subject to the volume $V(r, h) = \pi r^2 h = 355$.
The gradient equations are
$$4\pi r + 2\pi h = \lambda\,(2\pi rh), \qquad 2\pi r = \lambda\,(\pi r^2)$$
The second gives $\lambda = \frac{2}{r}$. Put that into the first:
$$4\pi r + 2\pi h = \frac{2}{r}(2\pi rh) = 4\pi h \quad\Longrightarrow\quad 4\pi r = 2\pi h \quad\Longrightarrow\quad h = 2r$$
Look at what happened. Before using the number 355 at all, we learned the
\emph{shape} of the best can: its height equals its diameter. That holds for a can of
any size. Then $\pi r^2(2r) = 355$ gives $r \approx 3.84$ cm and $h \approx 7.67$ cm.

(Real soda cans are about 12 cm tall and 6.6 cm wide, much taller than the ideal.
They are easier to hold, the top and bottom are made of thicker metal than the
sides, and a narrow can fits more neatly in a vending machine.)

\section{A Probability Example: The Most Likely Spinner}\index{maximum likelihood}
A spinner has three colors, but you do not know how big each section is. You spin it
20 times and get red 10 times, green 6 times, and blue 4 times. What are the most
believable probabilities $p$, $q$, and $s$ for red, green, and blue?

Spins are independent, so the probability of getting this particular sequence of
results is the product of the individual chances:
$$P = p^{10} q^{6} s^{4}$$
(For example, two heads in a row with a fair coin has probability
$\frac{1}{2} \cdot \frac{1}{2} = \frac{1}{4}$.) The probabilities must add to 1, so we
want to
$$\text{maximize } p^{10}q^{6}s^{4} \quad\text{subject to}\quad p + q + s = 1$$
This is called finding the \newterm{maximum likelihood} estimate: the probabilities that
make what you actually saw most likely.

Products of powers are awkward to differentiate, so take the natural log first. Since
$\ln$ always increases, $P$ and $\ln P$ are largest at the same point, and
$$\ln P = 10\ln p + 6\ln q + 4\ln s$$
The method works the same way with three variables: set
$\nabla(\ln P) = \lambda\nabla(p + q + s)$:
$$\frac{10}{p} = \lambda, \qquad \frac{6}{q} = \lambda, \qquad \frac{4}{s} = \lambda$$
So $p = \frac{10}{\lambda}$, $q = \frac{6}{\lambda}$, and $s = \frac{4}{\lambda}$. The constraint
gives $\frac{20}{\lambda} = 1$, so $\lambda = 20$, and
$$p = \frac{10}{20} = 0.5, \qquad q = \frac{6}{20} = 0.3, \qquad s = \frac{4}{20} = 0.2$$
The most likely probabilities are just the fractions you observed. That may sound
obvious, but now you have a proof, and the same argument works for any number of
outcomes and any counts. It is the starting point for much of statistics and machine
learning, which you will meet in later volumes.

\section{What $\lambda$ Means}\index{Lagrange multipliers!meaning of lambda}
The multiplier is more than a bookkeeping device. It tells you how much the best value
of $f$ changes when you loosen the constraint: if $c$ increases by a small amount
$\Delta c$, the optimum changes by about $\lambda\,\Delta c$.

Check it on the fence. With $x + y = c$, the best rectangle is a square with side
$\frac{c}{2}$, so the best area is $\frac{c^2}{4}$. Its derivative is $\frac{c}{2}$, which is 4
when $c = 8$: exactly the $\lambda$ we found. So one more meter of $x + y$ (two more
meters of fence) buys about 4 more square meters of garden. Economists call $\lambda$ a
\newterm{shadow price}: what one more unit of a limited resource is worth.

\section{Lagrange Multipliers with SymPy}
The equations can get messy, and a computer algebra system can solve them for you. Here
is the circle example in Python with the SymPy library:
\begin{minted}{python}
import sympy as sp

x, y, lam = sp.symbols("x y lambda", real=True)
f = x**2 + 2*y**2          # what we optimize
g = x**2 + y**2            # the constraint is g = 1

# grad f = lambda * grad g, together with the constraint
equations = [
    sp.Eq(sp.diff(f, x), lam * sp.diff(g, x)),
    sp.Eq(sp.diff(f, y), lam * sp.diff(g, y)),
    sp.Eq(g, 1),
]
for sol in sp.solve(equations, [x, y, lam], dict=True):
    print(sol, " f =", f.subs(sol))
\end{minted}
\begin{verbatim}
{lambda: 1, x: -1, y: 0}  f = 1
{lambda: 1, x: 1, y: 0}  f = 1
{lambda: 2, x: 0, y: -1}  f = 2
{lambda: 2, x: 0, y: 1}  f = 2
\end{verbatim}
All four candidate points appear, with their values of $f$.

\begin{Exercise}[title={On a Circle}, label=lm_circle]
Find the maximum and minimum of $f(x, y) = x + y$ on the circle $x^2 + y^2 = 8$. Draw a
sketch like figure \ref{fig:lagrange-contours} to check your answer.
\end{Exercise}
\begin{Answer}[ref=lm_circle]
$\nabla f = (1, 1)$ and $\nabla g = (2x, 2y)$, so $1 = 2\lambda x$ and $1 = 2\lambda y$. Thus
$x = y$, and $2x^2 = 8$ gives $x = y = \pm 2$. The maximum is $f(2, 2) = 4$ and the minimum
is $f(-2, -2) = -4$. In a sketch, the level lines $x + y = c$ are parallel lines, and the
two extreme ones just touch the circle at $(2, 2)$ and $(-2, -2)$.
\end{Answer}

\begin{Exercise}[title={Fencing by a River}, label=lm_river]
A farmer has 100 m of fence to enclose a rectangular field along a straight river. No
fence is needed along the river. Let $x$ be the side parallel to the river and $y$ each
side perpendicular to it.
\begin{enumerate}
\item Use Lagrange multipliers to find the largest possible area.
\item What is $\lambda$, and what does it tell the farmer?
\end{enumerate}
\end{Exercise}
\begin{Answer}[ref=lm_river]
\begin{enumerate}
\item Maximize $xy$ subject to $x + 2y = 100$. The equations are $y = \lambda$ and
$x = 2\lambda$, so $x = 2y$. Then $4y = 100$, so $y = 25$ m, $x = 50$ m, and the area is
1250 m$^2$.
\item $\lambda = 25$. Each extra meter of fence adds about 25 m$^2$ of field. (Check: with
$F$ meters of fence, the best area is $\frac{F^2}{8}$, and its derivative at $F = 100$ is
$\frac{100}{4} = 25$.)
\end{enumerate}
\end{Answer}

\begin{Exercise}[title={The Biggest Box}, label=lm_box]
A closed box has surface area $2(xy + yz + xz) = 24$. Find the largest possible volume
$xyz$. (Hint: with three variables, there are four equations.)
\end{Exercise}
\begin{Answer}[ref=lm_box]
The equations are $yz = 2\lambda(y + z)$, $xz = 2\lambda(x + z)$, and $xy = 2\lambda(x + y)$.
Multiply the first by $x$ and the second by $y$: $xyz = 2\lambda(xy + xz)$ and
$xyz = 2\lambda(xy + yz)$, so $xz = yz$, and $x = y$. In the same way, $y = z$. So the box is
a cube, $6x^2 = 24$ gives $x = 2$, and the largest volume is $8$.
\end{Answer}

\begin{Exercise}[title={An Ellipse}, label=lm_ellipse]
Find the maximum and minimum of $f(x, y) = xy$ on the ellipse $x^2 + 4y^2 = 8$. Then
check your answer with SymPy.
\end{Exercise}
\begin{Answer}[ref=lm_ellipse]
$y = 2\lambda x$ and $x = 8\lambda y$. Substituting, $y = 16\lambda^2 y$, so $\lambda = \pm\frac{1}{4}$
(the case $y = 0$ forces $x = 0$, which is not on the ellipse). With $\lambda = \frac{1}{4}$,
$x = 2y$, and $4y^2 + 4y^2 = 8$ gives $y = \pm 1$: the points $(2, 1)$ and $(-2, -1)$, where
$f = 2$. With $\lambda = -\frac{1}{4}$, $x = -2y$: the points $(-2, 1)$ and $(2, -1)$, where
$f = -2$. The maximum is 2 and the minimum is $-2$. In the SymPy program, use
\mintinline{python}{f = x*y} and \mintinline{python}{g = x**2 + 4*y**2}, with the
constraint \mintinline{python}{sp.Eq(g, 8)}.
\end{Answer}

\begin{Exercise}[title={The Most Uncertain Spinner}, label=lm_entropy]
A measure of how unpredictable a spinner with probabilities $p$, $q$, and $s$ is, called
its \newterm{entropy}, is
$$H = -p\ln p - q\ln q - s\ln s$$
Use Lagrange multipliers to find the probabilities that make $H$ as large as possible,
subject to $p + q + s = 1$. Does the answer make sense?
\end{Exercise}
\begin{Answer}[ref=lm_entropy]
$\frac{\partial H}{\partial p} = -\ln p - 1$, and similarly for $q$ and $s$. So
$-\ln p - 1 = \lambda$, $-\ln q - 1 = \lambda$, and $-\ln s - 1 = \lambda$, which means
$\ln p = \ln q = \ln s$, so $p = q = s = \frac{1}{3}$. The most unpredictable spinner gives
every color the same chance, which makes sense: any imbalance would make one color a
safer guess.
\end{Answer}
